%Chargement des paquets
\usepackage{amsmath}
\usepackage{amsfonts}
\usepackage{amssymb}
\usepackage{enumerate}
\usepackage{mathtools}
\usepackage{bbm}
\usepackage{xparse, etoolbox}
\usepackage{enumerate}
\usepackage{enumitem}
\usepackage{mathabx}
\usepackage{babel}[french]

%environment
\usepackage{amsthm}
\usepackage{keytheorems}
\usepackage{hyperref} 

%Interface théorème
\newkeytheorem{lem}[
	name = Lemme
]

\newkeytheorem{prop}[
	name = Proposition
]

\newkeytheorem{thm}[
	name = Théorème
]

\newkeytheorem{coro}[
	name = Corollaire
]

\renewcommand*{\proofname}{Démonstration}

\theoremstyle{definition}
\newtheorem{defn}{Définition}

\theoremstyle{definition}
\newtheorem*{nota}{Notation}

\theoremstyle{definition}
\newtheorem*{limi}{Liminaire}

\theoremstyle{remark}
\newtheorem{rmq}{Remarque}

\theoremstyle{plain}
\newtheorem*{fait}{Fait}


%Ensembles de nombres
\newcommand{\N} {\mathbb{N}}
\newcommand{\Ne}{\N^\ast}
\newcommand{\Z} {\mathbb{Z}}
\newcommand{\Q} {\mathbb{Q}}
\newcommand{\R} {\mathbb{R}}
\newcommand{\Rb}{\overline{\mathbb{R}}}
\newcommand{\Rp}{\R_+}
\newcommand{\Rm}{\R_-}
\newcommand{\K} {\mathbb{K}}
\newcommand{\Ket} {\mathbb{K^{\ast}}}
\newcommand{\Cx}{\mathbb{C}}

%Ensembles, tribus
\newcommand{\A}{\mathbb{A}}
\renewcommand{\P}{\mathcal{P}}
\newcommand{\T}{\mathcal{T}}
\newcommand{\F}{\mathcal{F}}
\newcommand{\C} {\mathcal{C}}
\newcommand{\ouv}{\mathcal{O}}
\newcommand{\B}{\mathcal{B}}
\newcommand{\M}{\mathcal{M}}
\newcommand{\etag}{\mathcal{E}}
\newcommand{\etagOp}{\etag^{+}(\Omega)}
\newcommand{\etagO}{\etag(\Omega)}
%%Lp avant quotientage
\newcommand{\Lic}{\mathcal{L}^1}
\newcommand{\Lpc}{\mathcal{L}^p}
\newcommand{\Linfc}{\mathcal{L}^{\infty}}
\newcommand{\Lifull}{\Lic(\Omega, \B, m)}
\newcommand{\Lpfull}{\Lpc(\Omega, \B, m)}
\newcommand{\Linffull}{\Linfc(\Omega, \B, m)}
%%Lp après quotientage
\newcommand{\Li}{L^1}
\newcommand{\Ld}{L^2}
\newcommand{\Lp}{L^p}
\newcommand{\Linf}{L^{\infty}}
\newcommand{\Listd}{\Li(\Omega, m)} 
\newcommand{\Ldstd}{\Ld(\Omega, m)} 
\newcommand{\Lpstd}{\Lp(\Omega, m)} 

%Quelques opérateurs
\DeclareMathOperator{\codim}{codim}
\DeclareMathOperator{\rg}{rg}
\DeclareMathOperator{\tr}{tr}
\DeclareMathOperator{\vect}{Vect}
\DeclareMathOperator{\ima}{Im}
\DeclareMathOperator{\id}{Id}
\DeclareMathOperator{\sign}{sign}
\DeclareMathOperator{\ext}{ext}
\newcommand{\ind}{\mathbbm{1}}
\newcommand*{\spr}[2]{\left(\,#1\,\middle|\,#2\,\right)}
\newcommand{\interior}[1]{%
  {\kern0pt#1}^{\mathrm{o}}%
}
\DeclareMathOperator{\card}{card}
\DeclareMathOperator{\ppcm}{ppcm}

%Commandes de pure flemme
\newcommand{\veps}{\varepsilon}
\newcommand{\intab}{\int_{a}^{b}}
\newcommand{\intR}{\int_{-\infty}^{\infty}}
\newcommand{\intinf}{\int_{- \infty}^{+ \infty}}
\newcommand{\equi}{\Leftrightarrow}
\newcommand{\Kev}{$\K$-espace vectoriel }
\newcommand{\Kevs}{$\K$-espaces vectoriels }
\newcommand{\Rev}{$\R$-espace vectoriel }
\newcommand{\Revs}{$\R$-espaces vectoriels }
\newcommand{\Cev}{$\Cx$-espace vectoriel }
\newcommand{\Cevs}{$\Cx$-espaces vectoriels }
\newcommand{\evn}{(E, \norm{.})}
\newcommand{\interf}[2]{[#1,#2]}
\newcommand{\limb}{\lim\limits_}
\newcommand{\lebn}{\lambda_n}
\newcommand{\pp}{\text{~presque partout}}
\newcommand{\derivp}[2]{\frac{\partial #1}{\partial #2}}
\newcommand{\famiu}{(u_i)_{i \in I}}
\newcommand{\sev}{\text{~s.e.v.}}
\newcommand{\Mnp}{M_{n,p}(\K)}
\newcommand{\Mn}{M_{n}(\K)}
\newcommand{\tq}{\text{~t.q.~}}
\newcommand{\mq}{~montrer~que~}
\newcommand{\et}{\quad \text{et} \quad}
\newcommand{\ou}{\text{~ou~}}
\newcommand{\Kx}{\K[X]}


%Commandes de gars chelou
\renewcommand{\subset}{\subseteq}

%Commande restriction, ty duckduckgo
\def\restr#1#2{\mathchoice
              {\setbox1\hbox{${\displaystyle #1}_{\scriptstyle #2}$}
              \restraux{#1}{#2}}
              {\setbox1\hbox{${\textstyle #1}_{\scriptstyle #2}$}
              \restraux{#1}{#2}}
              {\setbox1\hbox{${\scriptstyle #1}_{\scriptscriptstyle #2}$}
              \restraux{#1}{#2}}
              {\setbox1\hbox{${\scriptscriptstyle #1}_{\scriptscriptstyle #2}$}
              \restraux{#1}{#2}}}
\def\restraux#1#2{{#1\,\smash{\vrule height .8\ht1 depth .85\dp1}}_{\,#2}} 

%Quelques normes
\DeclarePairedDelimiter\abs{\lvert}{\rvert}
\DeclarePairedDelimiter\labs{\bigl\lvert}{\bigr\rvert}
\DeclarePairedDelimiter\norm{\lVert}{\rVert}
\DeclarePairedDelimiterXPP\normi[1]{}\lVert\rVert{_1}{\ifblank{#1}{\:\cdot\:}{#1}}
\DeclarePairedDelimiterXPP\normd[1]{}\lVert\rVert{_2}{\ifblank{#1}{\:\cdot\:}{#1}}
\DeclarePairedDelimiterXPP\normp[1]{}\lVert\rVert{_p}{\ifblank{#1}{\:\cdot\:}{#1}}
\DeclarePairedDelimiterXPP\normq[1]{}\lVert\rVert{_q}{\ifblank{#1}{\:\cdot\:}{#1}}
\DeclarePairedDelimiterXPP\norminf[1]{}\lVert\rVert{_\infty}{\ifblank{#1}{\:\cdot\:}{#1}}

%Bracket en angle
\DeclarePairedDelimiter\angl{\langle}{\rangle}

%Set, parenthèses
\newcommand{\set}[1]{\{ #1 \}}
\newcommand{\bigset}[1]{\bigl\{ #1 \bigr\}}
\newcommand{\Bigset}[1]{\Bigl\{ #1 \Bigr\}}
\newcommand{\bigpar}[1]{\bigl( #1 \bigr)}
\newcommand{\Bigpar}[1]{\Bigl( #1 \Bigr)}

%Opitmisation des opérateurs de comparaison
\DeclarePairedDelimiter\tio{o (}{)}
\DeclarePairedDelimiter\gro{O (}{)}


\pagestyle{fancy}

\fancyhead[C]{Critères de compacité}
\fancyhead[R]{}

\begin{document}

\underline{Source :} Rombaldi - Analyse - page 43 puis 47

\underline{Leçons :}
\begin{itemize}
\item Leçon 203 - Utilisation de la notion de compacité.
\end{itemize}

\begin{intro}
La notion de compacité présente une importante diversité de présentation parmi les auteurs des manuels les plus utilisés.
On constate par exemple que :
\begin{itemize}
\item dans son ouvrage de topologie générale, l'auteur Nawfal El Hage Hassan définit la compacité comme la propriété d'un espace
topologique séparé d'admettre un sous-recouvrement fini de tout recouvrement par des ouverts ;
\item dans son ouvrage destiné aux CPGE, Xavier Gourdon donne une définition similaire de la compacité, en se restreignant néanmoins
aux espaces métriques (donc séparés).
\item dans son ouvrage de préparation à l’agrégation, Jean-Etienne Rombaldi définit la compacité d'une partie d'un espace métrique
comme la propriété de toute suite de cette partie d'admettre une sous-suite convergente.
\end{itemize}
On voit ainsi se dessiner une équivalence entre deux propriétés et on comprend cette équivalence à gros traits : 
en un sens on peut recouvrir l'espace d'un nombre fini d'ensembles arbitrairement \og petits \fg{}, et toute suite aura nécessairement 
une infinité de points dans un de ces ensembles. 
Dans un second sens, tout recouvrement par des ouverts peut être assimilé à une suite de voisinage des points d'une suite. 
Or cette suite aura une infinité de points dans un ensemble arbitrairement \og petit \fg{}, et en construisant des ouverts
suffisamment \og grands \fg{} autour du nombre (fini) de points restants, on aura un recouvrement fini. \\
C'est en partant de cette intuition que l'on va démontrer l'équivalence les propriétés de Borel-Lebesgue et de Bolzano-Weierstrass
dans le cadre des espaces métriques.
\end{intro}

\begin{nota}
$(E,d)$ est un espace métrique et $K$ est une partie non vide de $E$.
\end{nota}

\begin{prop}[de Bolzano-Weierstrass]
De toute suite de points de $K$ on peut extraire une sous-suite qui converge dans $K$.
\end{prop}

\begin{lem}
\label{lem:1}
Si $K$ vérifie la propriété de Bolzano-Weierstrass, alors, pour tout rayon $r > 0$, 
il existe une suite finie $(a_k)_{1 \leq k \leq r}$ d'éléments de $K$ telle que :
\[ K \subset \bigcup_{k=1}^r B(a_k, r) \]
\end{lem}

\begin{proof}
Supposons qu'il existe $r > 0$ tel que cette propriété ne soit pas vraie. 
Prenons $u_0 \in K$, il existe alors $u_1 \in K \setminus B(u_0, r)$ et donc $d(u_0, u_1) \geq r$.
Supposons construite une suite finie $(u_0, \dots, u_n)$ de points de $K$ telle que $d(u_i, u_j) \geq r$ pour $i \neq j$.
Il existe alors $u_{n+1} \in K$ telle que $d(u_{n+1}, u_i) \geq r$ pour tout $i \in \set{1, \dots, n}$. 
Ainsi, on peut construire une suite $(u_n)_\N$ de $\K$ de laquelle il est impossible d'extraire une sous-suite convergente, ce qui est
contradictoire.
\end{proof}

\begin{lem}
\label{lem:2}
Soit $K$ vérifiant la propriété de Bolzano-Weierstrass et $(O_i)_I$ une famille d'ouverts recouvrant $K$.
Alors, il existe un réel $r>0$ tel que toute boule ouverte de centre $x \in K$ et de rayon $r$ soit contenue dans l'un des $O_i$.
\end{lem}

\begin{proof}
Supposons que pour tout réel $r>0$, il existe une boule ouverte de centre $x \in K$ et de rayon $r$ qui ne soit contenu dans aucun
des $O_i$. On construit alors une suite $B \left ( x_n, \dfrac1{n} \right )$ de boules ouvertes de centre $x_n \in K$ qui ne sont incluses
dans aucun des $O_i$. De la suite $(x_n)_\N$ on peut extraire sous-suite $(x_{\varphi(n)})_\N$ qui converge vers $x \in K$.
Il existe $i \in I$ tel que $x \in O_i$ et donc un rayon $r>0$ tel que $B(x, r) \subset O_i$.
En prenant $N \in \N$ tel que :
\[ \forall n \geq N, \quad d(x_{\varphi(n)}, x) < \dfrac{r}2 \et \abs*{\dfrac1{\varphi(n)}} < \dfrac{r}2, \]
on a alors, toujours pour $n \geq N$ :
\[ B \left ( x_{\varphi(n)}, \dfrac1{\varphi(n)} \right ) \subset B(x, r) \subset O_i , \]
ce qui est contradictoire.
\end{proof}

\begin{prop}[de Borel-Lebesgue]
De tout recouvrement de $K$ par des ouverts on peut extraire un sous-recouvrement fini.
\end{prop}

\begin{thm}
Les propriétés de Bolzano-Weierstrass et Borel-Lebesgue sont équivalentes.
\end{thm}

\begin{proof}
Soit $K$ vérifiant la propriété de Bolzano-Weierstrass, et $(O_i)_I$ une famille d'ouverts recouvrant $K$.
D'après le lemme $\ref{lem:2}$ il existe un rayon $r$ tel que toute boule ouverte de centre $x \in K$ de rayon $r$ soit contenu dans l'un
des $O_i$. Et d'après le lemme $\ref{lem:2}$, il existe une suite finie $(a_1, \dots, a_r)$ d'éléments K telle que :
\[  K \subset \bigcup_{k=1}^r B(a_k, r) \subset \bigcup_{k=1}^r O_{i_k} , \]
avec $O_{i_k}$ l'ouvert contenant $B(a_k, r)$. \\

Soit maintenant $K$ vérifiant la propriété de Borel-Legesgue et supposons qu'il existe une suite $(u_n)_\N$ de points de $K$ et
n'admettant pas de valeurs d'adhérence dans $K$. 
Donc, pour tout point $x \in K$, il existe un réel $r_x$ tel que $B(x, r_x)$ ne contienne qu'un nombre fini d'éléments de la suite.
On extrait un sous-recouvrement fini du recouvrement $\bigcup_{x \in K} B(x, r_x)$ de $K$ et on déduit que la suite $(u_n)_\N$ prend
un nombre fini de valeurs, donc admet une valeur d'adhérence, ce qui est contradictoire.
\end{proof}

On peut ainsi définir la notion de compacité :

\begin{defn}
Une partie $K$ d'un espace métrique $(E,d)$ est dite compacte si elle vérifie la propriété de Bolzano-Weierstrass ou la propriété de 
Borel-Lebesgue.
\end{defn}

\begin{rmq}
La propriété du lemme $\ref{lem:1}$ est appelée \og pré-compacité \fg{}. 
Une partie $K$ de $E$ est ainsi dite pré-compacte si elle vérifie cette propriété.
D'après ce que l'on a démontré, il est clair que la compacité implique la pré-compacité. \\
De la définition de la compacité, on tire les corollaires ci-dessous.
\end{rmq}

\begin{coro}
Une partie compacte de $E$ est fermée et bornée.
\end{coro}

\begin{proof}
Si $K$ est compact il vérifie la propriété de Bolzano-Weirstrass, donc est fermé. 
Il vérifie aussi la propriété de Borel-Lebesgue, donc deux points quelconques de $K$ seront distant d'au plus deux fois le maximum
des diamètres des ouverts recouvrant $K$ additioné du maximum des distances entre les ouverts du recouvrement considérés deux à
deux. 
\end{proof}

\begin{coro}
Une partie compacte de $E$ est complète.
\end{coro}

\begin{proof}
Soit $K$ un compact de $E$ et $(u_n)_\N$ une suite de points de $K$ de Cauchy. Cette suite admet une sous-suite $(u_{\varphi(n)})_\N$
convergente dans $K$ (ce dernier étant fermé), donc une valeur d'adhérence $l \in \K$.
La suite étant de Cauchy, pour tout $\veps >0$, on peut trouver $N \in \N$ tel que : 
\[ \forall p, q \geq N, \quad d(u_p, u_q) < \veps/2 . \]
Il existe aussi $N_1 \in \N$ tel que $\forall n \geq N, d(u_{\varphi(n)}, l) < \veps/2$.
Ainsi, en prenant $N_2 \in \N$ tel que $n \geq N_2 \Rightarrow \varphi(n) \geq \max(N, N_1)$, on a :
\[ \forall n \geq \max(N, N_2), \quad d(u_n, l) < d(u_n, u_{\varphi(n)}) + d(u_{\varphi(n)}, l) < \veps , \]
ce qui montre que la suite est convergente.
\end{proof}

La réciproque est évidemment fausse (considérer $\R$), on a néanmoins la propriété suivante :

\begin{thm}
Une partie complète et précompacte de $E$ est compacte.
\end{thm}

\begin{proof}
Supposons $K \subset E$ telle qu'énoncée et prenons $(u_n)_\N$ une suite d'éléments de $K$.
Par précompacité, on recouvre $K$ de boules ouvertes de rayon $1/2$, il existe alors une partie infinie de $\N$, que l'on note $J_0$,
et telle que la suite extraite $(u_n)_{J_0}$ soit incluse dans la boule $B(x_0, 1/2)$ avec $x_0 \in K$. \\
De même, on recouvre $K$ de boules ouvertes de rayon $1/4$, il existe alors une partie infinie de $J_0$, que l'on note $J_1$,
et telle que la suite extraite $(u_n)_{J_1}$ soit incluse dans la boule $B(x_1, 1/4)$ avec $x_0 \in K$. \\
En supposant ainsi construis les points $(x_0, \dots, x_n)$ et les parties $(J_0, \dots, J_n)$, on peut de même construire
$J_{n+1} \subset J_n$ infinie et telle que la suite $(u_n)_{J_{n+1}}$ soit incluse dans la boule $B(x_1, 1/2^{n+1})$ avec $x_{n+1} \in K$. 
On construit ainsi une suite décroissante (au sens de l'inclusion) de parties infinies de $\N$, 
on peut maintenant construire une extractrice $\varphi$ en posant :
\[ \varphi(0) = \min{J_0} \et \varphi(n+1) = \min( J_{n+1} \cap [\varphi(n) + 1 , \infty [ . \]
On note que $\varphi$ est bien définie car, pour tout $n \in \N, J_n$ n'est pas majoré (sinon il serait fini).
La suite extraite $(u_{\varphi(n)})_\N$ est de Cauchy, donc convergente, ce qui démontre que $K$ est compact.
\end{proof}

\end{document}